% Figure: pressure coefficient distribution over a transonic airfoil
% upper surface, showing the supersonic pocket terminated by a shock.
\begin{figure}[htbp]
\centering
\begin{tikzpicture}
\begin{axis}[
  width=12cm, height=8cm,
  xlabel={chordwise position $x/c$},
  ylabel={$-C_p$ (suction up)},
  xmin=0, xmax=1, ymin=-0.8, ymax=1.6,
  axis lines=left,
  y dir=normal,
  grid=major, grid style={enggray!20},
  tick label style={font=\scriptsize},
  label style={font=\small},
  legend style={font=\scriptsize, at={(0.97,0.97)}, anchor=north east},
  legend cell align=left,
]
% Cp* line: pressure coefficient at which local M=1 (sonic), here say -Cp*=0.7
\draw[enggray, dashed] (axis cs:0,0.7) -- (axis cs:1,0.7);
\node[enggray, font=\scriptsize, anchor=south east] at (axis cs:1,0.7) {$C_p^{*}$ (local $M=1$)};
% upper surface: strong leading-edge suction peak, supersonic plateau above Cp*, shock jump down
\addplot[engblue, very thick] coordinates {
  (0,-0.2) (0.03,1.2) (0.08,1.05) (0.15,0.95) (0.30,0.92) (0.45,0.90)
  (0.55,0.92)};
\addplot[engblue, very thick] coordinates {(0.55,0.92) (0.555,0.20)};
\addplot[engblue, very thick] coordinates {
  (0.555,0.20) (0.70,0.05) (0.85,-0.10) (1.0,-0.05)};
\addlegendentry{upper surface}
% lower surface: mild
\addplot[engred, very thick] coordinates {
  (0,-0.2) (0.05,-0.35) (0.2,-0.25) (0.5,-0.10) (0.8,0.0) (1.0,-0.05)};
\addlegendentry{lower surface}
% supersonic pocket label
\node[engblack, font=\scriptsize] at (axis cs:0.30,1.15) {supersonic pocket};
\node[engblack, font=\scriptsize, anchor=west] at (axis cs:0.57,0.55) {shock};
\end{axis}
\end{tikzpicture}
\caption[Transonic airfoil pressure distribution]{Pressure coefficient
over a transonic airfoil at a freestream Mach number above the
critical value, plotted as $-C_p$ so that suction is upward. The upper
surface accelerates the flow past the sonic line $C_p^{*}$ into a
local supersonic pocket; the pocket is closed by a near-normal shock
(the sharp recompression near mid-chord), downstream of which the
flow is subsonic again. The shock raises the pressure abruptly and,
if strong enough, separates the boundary layer behind it, the
mechanism of transonic buffet and drag rise. The lower surface stays
subsonic. The area enclosed between the upper and lower curves is
proportional to the section lift; the shock recompression is the
visible signature that the section is operating in the transonic
regime.}
\label{fig:vol03ch13:airfoil-cp}
\end{figure}
